\documentclass[10pt,a4paper]{amsart}
\usepackage[margin=19mm]{geometry}
\usepackage{amsmath,amssymb,mathtools}
\usepackage[scr=rsfs,cal=euler]{mathalfa}
\usepackage{xfrac}
\usepackage[unicode,hidelinks,bookmarks=false]{hyperref}
\newcommand{\ie}{i.e.\ }
\newcommand{\cf}{cf.\ }
\newcommand{\resp}{respectively\ }
\newcommand{\Subsection}{\subsection}
\providecommand{\nobreakdash}{\nobreak\hbox{-}}
\setlength{\emergencystretch}{2em}
\multlinegap0pt
\newtheorem{theorem}{Theorem}[section]
\newtheorem{lemma}[theorem]{Lemma}
\theoremstyle{remark}
\newtheorem{rema}[theorem]{Remark}
\providecommand{\eg}{e.g.\ }
\providecommand{\etc}{etc.\ }
\DeclareMathOperator{\Var}{Var}
\DeclareMathOperator{\ord}{ord}
\DeclareMathOperator{\lcm}{lcm}
\DeclareMathOperator{\fix}{fix}
\DeclareMathOperator{\Poisson}{Poisson}
\DeclareMathOperator{\expe}{e}

\newcommand\calC{\mathcal{C}}
\newcommand\ordn{\ord S_n}
\newcommand{\eps}{\epsilon}
\renewcommand\P{\mathbf{P}}
\newcommand\E{\mathbf{E}}
\newcommand\br[1]{\left({#1}\right)}
\newcommand\floor[1]{\left\lfloor{#1}\right\rfloor}
\newcommand\pfrac[2]{\br{\frac{#1}{#2}}}

\newcommand*{\mk}{\mkern -1mu}
\newcommand*{\Mk}{\mkern -2mu}
\newcommand*{\mK}{\mkern 1mu}
\newcommand*{\MK}{\mkern 2mu}

\title[Poisson approximation under cycle-type restrictions]{Random generation with cycle type restrictions: original Theorem 3.10}
\author{Sean Eberhard and Daniele Garzoni}
\date{Source: Algebraic Combinatorics 4(1) (2021), 1--25; DOI 10.5802/alco.149}
\begin{document}\maketitle
\noindent\textit{Editorial scope.} The counted unit is original Theorem 3.10. The complete original Section 3 bounding common small orbits and establishing the mixed factorial-moment estimates is uncounted source-local core. Subsequent imprimitive/primitive generation and fixed-order applications are omitted. The reference to the omitted Section 4 is a link to the original article; no author mathematical formula or assumption is altered. The retained passages contain no external bibliography citation.
\section{Original notation}
Let $\calC, \calC' \subset S_n$ be fixed conjugacy classes. For each $j$ let $c_j$ be the number of $j$-cycles in a representative element of $\calC$, and similarly $c'_j$ for $\calC'$. Let $\pi\in\calC, \pi'\in\calC'$ be chosen uniformly at random.

Much of our notation has already appeared. We consistently write $\pi$ and $\pi'$ for our two permutations, which we assume are random subject to having $c_j$ and $c'_j$ $j$-cycles, respectively, for each $j$. We will write $c$ and $c'$ for the total number of cycles:
\begin{align*}
 c&= c_1 + \cdots + c_n \\
 c'&= c'_1 + \cdots + c'_n.
\end{align*}
Occasionally, we will denote cycle types using exponent notation: Thus for example the conjugacy class $\calC$ consists of all permutations with cycle type $(1^{c_1}, 2^{c_2}, \dots, n^{c_n})$.

We write $\Omega$ for our ground set of size $n$, and we write $S_n$ and $A_n$ for the symmetric and alternating groups on $\Omega$. If $G$ acts on a set $X$, $g\in G$, $x \in X$, we write $x^g$ for the image of $x$ under $g$. In particular, if $g, h \in G$ we write $h^g$ for $g^{-1}hg$. Thus if $\fix_X(g)$ is the fixed point set of $g$ then we have
\[
 \fix_X(g^h) = \fix_X(g)^h.
\]

We use standard big-O and little-o without reservation, all understood with respect to the limit $n \to\infty$.
That is, if $x$ and $y$ are quantities depending on $n$, $x = O(y)$ means there is an constant $C$ such that $|x| \leq C y$ for all sufficiently large $n$,
while $x = o(y)$ means $x/y \to 0$ as $n\to\infty$.
Similarly, $x = \Omega(y)$ means there is a constant $c>0$ such that $x \geq c y$, and $x = \omega(y)$ means $x/y \to \infty$.

Additionally, for positive quantities we use the Vinogradov notation $x \ll y$ to mean $x = O(y)$, and $x \asymp y$ to mean $x \ll y$ and $x \gg y$ (\ie $x/y$ is bounded above and below by positive constants).
We will use subscripts such as $O_{k,m}$ or $\ll_{k,m}$ to warn that the implicit constant may depend on the parameters $k, m$.

\setcounter{section}{2}
\section{Intransitive subgroups} \label{sec:intransitive}

Let $G = \langle \pi, \pi'\rangle$. As so often the case in random generation problems in $S_n$, the main obstruction is transitivity: If $G$ is transitive then very likely $G \geq A_n$ (we will prove this in Section~\href{https://alco.centre-mersenne.org/articles/10.5802/alco.149/}{4}). In other words, our main job is to estimate the probability that $\pi$ and $\pi'$ have a common fixed set of some size $k$.

Let $N_k$ be the number of orbits of $G$ of size $k$, in other words the number of $k$-sets fixed by $G$ and on which $G$ acts transitively. Let
\[
 N = \sum_{1\leq k \leq n/2} N_k.
\]
Our goal is to estimate $\P(N = 0)$. We will accomplish this in two stages:
\begin{enumerate}
 \item\label{sect2.5_1} We will bound $\lambda = \E N$. This will already establish what we need for the almost-sure problem, since $\P(N = 0) \geq 1 - \lambda$.
 \item\label{sect2.5_2} We will prove a Poisson-type approximation for $N$ which shows that $\P(N=0) \approx \expe^{-\lambda}$.
\end{enumerate}

\subsection{Bounding \texorpdfstring{$\lambda = \E N$}{lambda = EN}}

There are many ways we can form a fixed $k$-set from the cycles of $\pi$ and $\pi'$. For any two partitions
\begin{align}
 k &= d_1 1 + d_2 2 + \cdots + d_k k, \label{eqn:d-partition}\\
 k &= d'_1 1 + d'_22 + \cdots + d'_k k, \label{eqn:d'-partition}
\end{align}
such that $d_j \leq c_j, d'_j \leq c'_j$ for each $j$, we can form a fixed $k$-set by taking $d_j$ $j$-cycles from $\pi$ and $d'_j$ $j$-cycles from $\pi'$, for each $j$. The probability that these $k$-sets coincide after $\pi'$ is conjugated by $\sigma$ is $1 / \binom{n}{k}$. Thus we have
\begin{equation} \label{eqn:ENk-formula}
 \E N_k = \sum_{\eqref{eqn:d-partition}, \eqref{eqn:d'-partition}} \frac{\binom{c_1}{d_1} \cdots \binom{c_k}{d_k} \binom{c'_1}{d'_1} \cdots \binom{c'_k}{d'_k}}{\binom{n}{k}} p(d_1, \dots, d_k ; d'_1, \dots, d'_k),
\end{equation}
where $p(d_1, \dots, d_k ; d'_1, \dots, d'_k)$ denotes the probability that random permutations $\tau, \tau' \in S_k$ with cycle types $(1^{d_1}, \dots, k^{d_k})$ and $(1^{d'_1}, \dots, k^{d'_k})$ generate a transitive subgroup.\footnote{Of course, this is the very thing we are trying to estimate in this section. This ``self-similarity'' is largely incidental: When we write $p(d_1, \dots, d_k ; d'_1, \dots, d'_k)$ we usually have in mind $k, d_1, \dots, d'_k$ of bounded size, whereas in the bigger picture of this section we are concerned with asymptotics.}

In this subsection we simply ignore $p(d_1, \dots, d_k ; d'_1, \dots, d'_k)$, bounding it by $1$. Thus from~\eqref{eqn:ENk-formula} we have simply
\[
 \E N_k \leq \frac{ f_k f'_k }{\binom{n}{k}},
\]
where
\begin{equation}\label{fk}
 f_k = \sum_{\eqref{eqn:d-partition}} \binom{c_1}{d_1} \cdots \binom{c_k}{d_k}
\end{equation}
is the number of $k$-sets fixed by $\pi$, and similarly $f'_k$. (In the next subsection we will need the full force of~\eqref{eqn:ENk-formula}.)

\begin{lemma} \label{lem:fkx-bound}
We have
\[
 f_k \leq \min_{x > 0}\, x^{-k} (1+x)^{c_1} (1+x^2)^{c_2} \cdots (1+x^k)^{c_k}.
\]
\end{lemma}
\begin{proof}
This follows from
\begin{align*}
 f_k x^k
 &= \sum_{\eqref{eqn:d-partition}} \binom{c_1}{d_1} \cdots \binom{c_k}{d_k} x^{d_1 1 + \cdots + d_k k} \\
 &\leq \sum_{d_1, \dots, d_k\geq 0} \binom{c_1}{d_1} \cdots \binom{c_k}{d_k} x^{d_1 1 + \cdots + d_k k} \\
 &= (1+x)^{c_1} (1+x^2)^{c_2} \cdots (1+x^k)^{c_k}.\qedhere
\end{align*}
\end{proof}

Define the function $h$ by
\[
 h(x) = x \log \frac1x + (1-x) \log\frac1{1-x}.
\]
It is well known that
\[
 k^{-1/2} \expe^{h(k/n)n} \ll \binom{n}{k} \leq \expe^{h(k/n)n} \qquad (1\leq k\leq n/2).
\]
The latter inequality follows from
\[
 \binom{n}{k} x^k \leq (1+x)^n,
\]
with $x = k/(n-k)$, since in this case
\begin{equation}
 \label{h-identity}
 x^{-k}(1+x)^n = \expe^{h(k/n)n}.
\end{equation}
The former inequality follows from Stirling's approximation.

\begin{lemma} \label{fk-bound}
For any $\eps \in (0, 1/2)$ we have
\[
 f_k \leq
 \pfrac{c_1}{k}^k \expe^{O(\eps^{-2}k)} +
 \expe^{h\pfrac{k}{2c_2} c_2 + \eps k} +
 \pfrac{n}{k}^{k/3} \expe^{O(\eps^{-2}k)}
 .
\]
\end{lemma}

To be clear, the assertion is that there is a constant $C$ not depending on $\eps$ such that, for all $\eps \in (0, 1/2)$,
\[
 f_k \leq
 \pfrac{c_1}{k}^k \expe^{C \eps^{-2}k} +
 \expe^{h\pfrac{k}{2c_2} c_2 + \eps k} +
 \pfrac{n}{k}^{k/3} \expe^{C \eps^{-2}k}
 .
\]
Hence $\eps$ is an optimizable parameter.

Note that the middle term here is of the form
\[
 \pfrac{c_2}{k}^{k/2} \expe^{O(k)}.
\]
The more precise bound is important to us, however.

\begin{proof}
First note that we may assume $k \leq 3c/2$, for if $k \geq 3c/2$ then directly from~\eqref{fk} we have
\[
 f_k \leq 2^c \leq 2^{2k/3},
\]
which is subsumed by the third term in the claim.

Now assume $k \leq 3c/2$. By the previous lemma we have, if $x \leq 1$,
\begin{equation} \label{eqn:fkx-bound}
 f_k \leq x^{-k} (1+x)^{c_1} (1+x^2)^{c_2} (1+x^3)^{c}.
\end{equation}
We consider three cases\footnote{To motivate the three cases, note that the unique minimizer of~\eqref{eqn:fkx-bound} satisfies
\[
-\frac{k}{x} + \frac{c_1}{1+x} + \frac{2c_2 x}{1+x^2} + \frac{3c x^2}{1+x^3} = 0,
\]
and hence
\[
 c_1 \pfrac{x}{1+x} + 2c_2 \pfrac{x^2}{1+x^2} + 3c \pfrac{x^3}{1+x^3} = k.
\]
The solution to this equation is comparable to the smallest of $k/c_1$, $(k/c_2)^{1/2}$, $(k/c)^{1/3}$.
We consider these three cases separately, taking extra care with the $c_2$ dependence.}
depending on which of
\[
 c_1/k,\quad \frac12 \eps (c_2 / k)^{1/2},\quad (c / k)^{1/3}
\]
is largest.
\begin{enumerate}
 \item\label{proof_lemm3.2_1} First suppose $c_1 / k \geq \frac12 \eps (c_2 / k)^{1/2}, (c / k)^{1/3}$. Take $x = k / c_1$ in~\eqref{eqn:fkx-bound}. The result is
 \[
 f_k \leq \pfrac{c_1}{k}^k \expe^{\pfrac{k}{c_1} c_1 + \pfrac{k}{c_1}^2 c_2 + \pfrac{k}{c_1}^3 c} \leq \pfrac{c_1}{k}^k \expe^{O(\eps^{-2} k)}.
 \]
 \item\label{proof_lemm3.2_2} Next suppose $\frac12 \eps(c_2 / k)^{1/2} \geq c_1/k, (c/k)^{1/3}$. Note the latter condition implies $k \leq \eps^6 c_2 \leq c_2 / 2$. Take
 \[
 x = \pfrac{k}{2c_2 - k}^{1/2}
 \]
 in~\eqref{eqn:fkx-bound}. Using
 \[
 f_k \leq x^{-k} (1+x^2)^{c_2} \expe^{x c_1 + x^3 c_3}
 = \br{(x^2)^{-k} (1+x^2)^{2c_2}}^{1/2} \expe^{x c_1 + x^3 c_3}
 \]
 and~\eqref{h-identity}, we have
 \[
 f_k
 \leq \expe^{h \pfrac{k}{2c_2} c_2 + \pfrac{k}{2c_2-k}^{1/2} c_1 + \pfrac{k}{2c_2-k}^{3/2} c}.
 \]
 Now by hypothesis
 \[
 \pfrac{k}{2c_2 - k}^{1/2} c_1 \leq \pfrac{k}{c_2}^{1/2} c_1 \leq \eps k / 2
 \]
 and
 \[
 \pfrac{k}{2c_2-k}^{3/2} c \leq \pfrac{k}{c_2}^{3/2} c \leq (\eps / 2)^3 k \leq \eps k / 2,
 \]
 so
 \[
 f_k \leq \expe^{h \pfrac{k}{2c_2} c_2 + \eps k}.
 \]
 \item\label{proof_lemm3.2_3} Finally, suppose $(c / k)^{1/3} \geq \frac12 \eps (c_2 / k)^{1/2}, c_1 / k$. Take $x = (k/c)^{1/3}$ in~\eqref{eqn:fkx-bound}. The result is
 \[
 f_k \leq \pfrac{c}{k}^{k/3} \expe^{\pfrac{k}{c}^{1/3} c_1 + \pfrac{k}{c}^{2/3} c_2 + \pfrac{k}{c} c} \leq \pfrac{c}{k}^{k/3} \expe^{O(\eps^{-2} k)}.
 \]
\end{enumerate}
Since $c \leq n$ this proves the lemma.
\end{proof}

\begin{lemma} \label{lem:Nk-bound-small-k}
We have
\[
 \E N_k \ll \left(
 \frac{(c_1 + c_2^{1/2}) (c'_1 + {c'_2}^{1/2})}{n} + \frac{c_1 + c'_1}{n^{2/3}} + \pfrac{k}{n}^{1/6}
 \right)^k \expe^{O(k)}.
\]
Moreover, if $c_2 + c'_2 = n - \Omega(n)$ then
\begin{align*}
 \E N_k
% &
 \Mk \ll \Mk \frac1{k^{k/3}}
 \left(
 \frac{ c_1 c'_1 \Mk +\Mk c_1 {c'_2}^{1/2} \Mk +\Mk {c_2}^{1/2} c'_1}{n} \Mk +\Mk \frac{c_1 \Mk +\Mk c'_1}{n^{2/3}}
 \right)^k O(1)^k %\\
% &\qquad 
 \Mk +\Mk (1 \Mk -\Mk \Omega(1))^k + O(k/n)^{k/6}.
\end{align*}
\end{lemma}
\begin{proof}
We have
\[
 \E N_k \leq f_k f'_k / \binom{n}{k}.
\]
Bounding each of $f_k$ and $f'_k$ using the previous lemma and expanding the product, we get
\begin{align*}
 \E N_k 
 &\ll \frac1{k^k} \pfrac{c_1 c'_1}{n}^k \expe^{O(\eps^{-2} k)} \\
 &\qquad + \frac1{k^{k/2}} \pfrac{c_1 {c_2'}^{1/2} + {c_2}^{1/2} c'_1}{n}^k \expe^{O(\eps^{-2} k)} \\
 &\qquad + \frac1{k^{k/3}} \left(\frac{c_1 + c'_1}{n^{2/3}} \right)^k \expe^{O(\eps^{-2} k)} \\
 &\qquad + \left(\frac{k}{n} \right)^{k/6} \expe^{O(\eps^{-2} k)} \\
 &\qquad + k^{1/2} \expe^{h\left(\frac{k}{2c_2}\right) c_2 + h\left(\frac{k}{2c_2'}\right)c'_2 - h(k/n)n + 2 \eps k}.
\end{align*}
Here we used the simple bound $\binom{n}{k} \geq (n/k)^k$ for every term except the one involving both $c_2$ and $c'_2$, where we used the more precise bound $\binom{n}{k} \gg k^{-1/2} \expe^{h(k/n) n}$. Additionally, for the other terms involving $c_2$ or $c'_2$ we used
\[
 \expe^{h\pfrac{k}{2c_2}c_2} \leq \pfrac{c_2}{k}^{k/2} \expe^{O(k)}.
\]
If we apply this bound also to the $c_2, c'_2$ term, and bound various things crudely, then we get the first statement in the lemma.

To get the second statement, the only term which requires further comment is the $c_2, c'_2$ term. Since $h$ is concave and $h'(x) = \log(1/x - 1)$, we have
\[
 h(x) \leq h(k/n) + (x - k/n) \log(n/k - 1).
\]
Hence, by a short calculation,
\begin{align*}
 h\pfrac{k}{2c_2} c_2 + h\pfrac{k}{2c'_2} c'_2 - h\pfrac{k}{n} n
 &\leq -(n - c_2 - c'_2) \log \pfrac{n}{n-k} \\
 &\leq -(n - c_2 - c'_2) k/n.
\end{align*}
Assuming $c_2 + c'_2 = n - \Omega(n)$, this is $-\Omega(k)$, so the lemma follows from taking $\eps$ appropriately small.
\end{proof}

The above lemma suffices as long as $k = o(n)$. For larger $k$ it is convenient to argue a little differently (and more simply).

\begin{lemma} \label{lem:Nk-bound-large-k}
Assume $c_1 + c'_1 = o(n)$ and $c+c' \leq n - \Omega(n)$. Then for $k\leq n/2$ we have
\[
 \E N_k \leq \expe^{-\Omega(k) + o(n)}.
\]
\end{lemma}
\begin{proof}
From Lemma~\ref{lem:fkx-bound} we have, for $0 < x < 1$,
\[
 f_k f'_k \leq x^{-2k} (1+x)^{c_1 + c'_1} (1+x^2)^{c + c'}.
\]
Take $x = (k/(n-k))^{1/2}$. We have
\[
 (1+x)^{c_1 + c'_1} \leq \expe^{x(c_1 + c'_1)} \leq \expe^{(2k/n)^{1/2} (c_1 + c'_1)}
\]
and
\[
 (1+x^2)^{c+c'} = (1+x^2)^n (1-k/n)^{n - c - c'} \leq (1+x^2)^n \expe^{-k(n - c- c')/n}.
\]
Hence
\[
 f_k f'_k \leq \expe^{h(k/n)n} \expe^{(2k/n)^{1/2}(c_1 + c'_1) - k(n-c-c')/n},
\]
and
\[
 \E N_k \leq \frac{f_k f'_k}{\binom{n}{k}} \leq \expe^{O(\log k) + (2k/n)^{1/2} (c_1 + c'_1) - k(n-c-c')/n}.
\]
Assuming $c_1+c'_1 = o(n)$ and $c + c' \leq (1-\eps) n$, this is bounded by
\[
 \expe^{O(\log k) + o((2k n)^{1/2}) - \eps k} = \expe^{o(n) - \eps k},
\]
as claimed.
\end{proof}

The following theorem follows by combining the previous two lemmas.

\begin{theorem} \label{thm:main-ENk-bound}
$\ $
\begin{enumerate}
 \item If $c_1 + c'_1 = o(n^{2/3})$ and $(c_1 + {c_2}^{1/2}) (c'_1 + {c'_2}^{1/2}) = o(n)$ then $\E N = o(1)$.
 \item If $c_1 + c'_1 = O(n^{2/3})$, $(c_1 + {c_2}^{1/2}) (c'_1 + {c'_2}^{1/2}) = O(n)$, and $c_2 + c'_2 = n - \Omega(n)$, then $\E N = O(1)$.
\end{enumerate}
\end{theorem}

\subsection{A Poisson approximation} \label{subsec:possion}

In the previous subsection we bounded the expectation $\E N$ of $N$. In this subsection we apply the method of moments to show that, under similar hypotheses, if $\E N$ is bounded then in fact
\[
 \P(N = 0) \approx \expe^{-\E N}.
\]
The method of moments depends on being able to prove moment estimates of the form
\[
 \E N(N-1) \cdots (N-m+1) \approx (\E N)^m.
\]
We therefore now turn our attention to the estimation of the moments and mixed moments of $(N_k)_{k\geq 1}$.

Recall from~\eqref{eqn:ENk-formula} that
\[
 \E N_k = \sum_{\eqref{eqn:d-partition}, \eqref{eqn:d'-partition}} \frac{ \prod_{j=1}^k \binom{c_j}{d_j} \binom{c'_j}{d'_j}}{\binom{n}{k}} p(d_1, \dots, d_k; d'_1, \dots, d'_k).
\]
We first ask to what extent we have, for $k = O(1)$,
\[
 \E N_k \approx \sum_{\eqref{eqn:d-partition}, \eqref{eqn:d'-partition}} \frac{ \prod_{j=1}^k \frac{c_j^{d_j}}{d_j!} \frac{{c'_j}^{d'_j}}{d'_j!}}{\frac{n^k}{k!}} p(d_1, \dots, d_k; d'_1, \dots, d'_k).
\]
We have
\[
 \binom{c}{d} = (1 + O_d(1/c)) \frac{c^d}{d!},
\]
so the approximation is appropriate for the factors in which $c_j, c'_j$ are large, as well as the ones in which $d_j \leq 1$. Our contention is that the other terms do not contribute significantly, so overall the approximation is fine.

\begin{lemma}
We have
\[
 f_k \ll_k (c_1 + c^{1/2})^k.
\]
Moreover, for any particular index $j$ and any $t\geq 0$ we have
\[
 \sum_{\eqref{eqn:d-partition}, d_j \geq t} \binom{c_1}{d_1} \cdots \binom{c_k}{d_k} \ll_k \pfrac{c_j}{(c_1 + c^{1/2})^j}^t (c_1 + c^{1/2})^k.
\]
\end{lemma}
\begin{proof}
We revisit the calculation of the previous section, now armed with the restrictive hypothesis $k = O(1)$. From~\eqref{fk},
\[
 f_k \asymp_k \max_\eqref{eqn:d-partition} c_1^{d_1} \cdots c_k^{d_k}.
\]
The maximum is not difficult to analyze: in fact the maximum over real $d_i$ is precisely
\[
 \max \{ c_1^k, c_2^{k/2}, \dots, c_k^{k/k}\}.
\]
We bound this crudely by
\[
 \max\{c_1^k, c^{k/2}\} \asymp_k (c_1 + c^{1/2})^k.
\]

Now consider the part of the sum in which $d_j \geq t$. By the same token we have
\begin{align*}
 \max_{\eqref{eqn:d-partition}, d_j \geq t} c_1^{d_1} \cdots c_k^{d_k}
 &= c_j^t \max_{d_1 1 + \cdots + d_k k = k - tj} c_1^{d_1} \cdots c_k^{d_k} \\
 &\ll_k c_j^t (c_1 + c^{1/2})^{k-tj}.
\end{align*}
This proves the lemma.
\end{proof}

\begin{lemma}
\begin{align*}
 \E N_k
 &= \sum_{\eqref{eqn:d-partition}, \eqref{eqn:d'-partition}} \frac{ \prod_{j=1}^k \frac{c_j^{d_j}}{d_j!} \frac{{c'_j}^{d'_j}}{d'_j!}}{\frac{n^k}{k!}} p(d_1, \dots, d_k; d'_1, \dots, d'_k) \\
 &\qquad + O_k\left( \frac1{\min\{c_1 + c^{1/2}, c'_1 + {c'}^{1/2}\}^{2/3}}
 \frac{((c_1 + c^{1/2})(c'_1 + {c'}^{1/2}))^k}{n^k}
 \right).
\end{align*}
\end{lemma}

\begin{proof}
Let $\delta$ be a parameter, and split the sum according to whether there is any $j$ such that $d_j \geq 2$ and $c_j \leq \delta (c_1 + c^{1/2})^j$, or $d'_j\geq 2$ and $c'_j \leq \delta (c'_1 + {c'}^{1/2})^j$. By the previous lemma, the part of the sum where there is some such $j$ is bounded by
\[
 O_k(\delta^2) (c_1 + c^{1/2})^k (c'_1 + {c'}^{1/2})^k.
\]
On the other hand, if $c_j \geq \delta (c_1 + c^{1/2})^j$ whenever $d_j \geq 2$ then we have
\[
 \binom{c_1}{d_1} \cdots \binom{c_k}{d_k}
 = \frac{c_1^{d_1}}{d_1!} \cdots \frac{c_k^{d_k}}{d_k!} \left(1 + O_k\pfrac{1}{\delta (c_1 + c^{1/2})} \right),
\]
\looseness-1
and similarly for the primed variables. Thus the error in our approximation is bounded by
\[
 O_k\left(\delta^2 + \frac1{\delta (c_1 + c^{1/2})} + \frac1{\delta (c'_1 + {c'}^{1/2})} \right)
 \frac{(c_1 + c^{1/2})^k (c'_1 + {c'}^{1/2})^k}{n^k}.
\]
To get the best bound we take
\[
 \delta^3 = \min\{c_1 + c^{1/2}, c'_1 + {c'}^{1/2}\}^{-1}.
\]
This gives the claimed bound.
\end{proof}

There is a similar estimate for the mixed moment
\[
 \E[(N_1)_{m_1} \cdots (N_k)_{m_k}].
\]
Here we write $(x)_{m_i} = x(x-1) \cdots (x-m_i+1)$ for the ``falling factorial'', with the standard convention $(x)_0=1$. Let
\begin{align*}
 m &= m_1 + \cdots + m_k, \\
 M &= m_1 1 + \cdots + m_k k.
\end{align*}
Define $k_1, \dots, k_m$ by
\[
 (k_1, \dots, k_m) = (
 \overbrace{1, \dots, 1}^{m_1},
 \dots,
 \overbrace{k, \dots, k}^{m_k}
 ),
\]
\ie $k_i = j$ if
\[
 m_1 + \cdots + m_{j-1} \leq i \leq m_1 + \cdots + m_j.
\]
For example, if $(m_1, m_2, m_3) = (3, 2, 1)$ then $m = 6$, $M = 10$, and
\[
 (k_1, k_2, k_3, k_4, k_5, k_6)
 = (1, 1, 1, 2, 2, 3).
\]
The product
\[
 (N_1)_{m_1} \cdots (N_k)_{m_k}
\]
counts $m$-tuples of distinct orbits of $\langle \pi, \pi'\rangle$, of which $m_1$ are $1$-sets, $m_2$ are $2$-sets, and so on, all these sets being \emph{disjoint} (this is why we count orbits rather than fixed sets). To choose such sets we have to choose, for every such size $k_i$, $d_{i1}$ $1$-cycles, $d_{i2}$ $2$-cycles, $\dots$ of $\pi$, and likewise $d'_{i1}$ $1$-cycles, $d'_{i2}$ $2$-cycles, $\dots$ of $\pi'$. This reasoning leads to the following estimate.

\begin{lemma}
\begin{align*}
 &\E[(N_1)_{m_1} \cdots (N_k)_{m_k}] \\
 &\qquad = \sum \prod_{i=1}^m \frac{\prod_{j=1}^{k} \frac{c_j^{d_{ij}}}{d_{ij}!}\frac{{c'_j}^{d'_{ij}}}{d'_{ij}!}}{\frac{n^{k_i}}{k_i!}} p(d_{i1}, \dots, d_{ik} ; d'_{i1}, \dots, d'_{ik}) \\
 &\qquad\qquad + O_M\left(
 \frac1{\min\{c_1 + c^{1/2}, c'_1 + {c'}^{1/2}\}^{2/3}}
 \frac{((c_1 + c^{1/2})(c'_1 + {c'}^{1/2}))^M}{n^M}
 \right).
\end{align*}
The sum runs over all $(d_{ij})$, $(d'_{ij})$ such that
\[
 d_{i1} 1 + \cdots + d_{ik} k =
 d'_{i1} 1 + \cdots + d'_{ik} k = k_i \qquad (i \in \{1, \dots, m\}).
\]
\end{lemma}

\begin{proof}
We have\footnote{We are using a slightly nonstandard, but backwards-compatible, notation for multinomial coefficients: $\binom{n}{a, b} = \frac{n!}{a!b!(n-a-b)!}$, \etc}
\begin{equation*}
% &
 \E[(N_1)_{m_1} \cdots (N_k)_{m_k}] %\\
% &\qquad 
 \Mk = \Mk \sum \Mk \frac{\prod_{j=1}^{k} \binom{c_j}{d_{1j}, \dots, d_{mj}} \binom{c'_j}{d'_{1j}, \dots, d'_{mj}}}
 {\binom{n}{k_1, \dots, k_m}}
 \mk \prod_{i=1}^m \mk p(d_{i1}, \dots, d_{ik} ; d'_{i1}, \dots, d'_{ik}).
\end{equation*}
The rest of the proof is the same as in the previous lemma.
\end{proof}

Write $N_{\leq k} = N_1 + \cdots + N_k$ and $N_{>k} = N - N_{\leq k}$.
These are the number of orbits of $G$ of size at most $k$ and of size greater than $k$, respectively.

\begin{theorem} \label{thm:poisson-approximation}
Let
\[
 B = \max \{(c_1 + c^{1/2}) (c'_1 + {c'}^{1/2}) / n, 1\}.
\]
Then
\[
 \E[(N_1)_{m_1} \cdots (N_k)_{m_k}] - (\E N_1)^{m_1} \cdots (\E N_k)^{m_k}
 \ll_M
 \frac{B^M}{\min\{c_1 + c^{1/2}, c'_1 + {c'}^{1/2}\}^{2/3}}.
\]
Thus also, for any $k,m$,
\[
\E[(N_{\leq k})_m] - (\E N_{\leq k})^m 
 \ll_{k,m}
 \frac{B^{km}}{\min\{c_1 + c^{1/2}, c'_1 + {c'}^{1/2}\}^{2/3}}
 .
\]
\end{theorem}

\begin{proof}
The first bound follows immediately from the previous two lemmas. For the second bound, we use
\[
 (N_{\leq k})_m = m! \sum_{m_1 + \cdots + m_k = m} \frac{(N_1)_{m_1}}{m_1!} \cdots \frac{(N_k)_{m_k}}{m_k!},
\]
and similarly
\[
 (\E N_{\leq k})^m = m! \sum_{m_1 + \cdots + m_k = m} \frac{(\E N_1)^{m_1}}{m_1!} \cdots \frac{(\E N_k)^{m_k}}{m_k!}.
\]
Note that the value of
$
 M = 1 m_1 + \cdots + k m_k
$
over
$
 m_1 + \cdots + m_k = m
$
ranges between $m$ and $km$. This proves the theorem.
\end{proof}


\begin{theorem}
Assume $c_1 + c'_1 = o(n^{2/3})$, $(c_1 + c_2^{1/2}) (c'_1 + {c_2'}^{1/2}) = O(n)$, and $c_2 + c'_2 = n - \Omega(n)$. Then 
\[
 \P(N = 0) = \expe^{-\E N} + o(1).
\]
\end{theorem}

\begin{proof}
By Lemmas~\ref{lem:Nk-bound-small-k} and~\ref{lem:Nk-bound-large-k} from the previous subsection, there is a $k = O_\eps(1)$ such that $\E N_{>k}$ is smaller than $\eps$, and
\[
 \P(N = 0) \leq \P(N_{\leq k} = 0) \leq \P(N = 0) + \P(N_{>k} > 0),
\]
so the main thing to understand is $\E N_{\leq k}$. From Bonferroni's inequalities we have, for any $M\geq 0$,
\[
 1_{N_{\leq k}=0} = \sum_{m=0}^{M-1} (-1)^m \binom{N_{\leq k}}{m} + O\left(\binom{N_{\leq k}}{M}\right).
\]
(The left-hand side of this equation is the indicator that $N_{\leq k} = 0$.)
Therefore, from the previous theorem,
\begin{align*}
 \P(N_{\leq k} = 0)
 &= \sum_{m=0}^{M-1} (-1)^m \frac{ \E [(N_{\leq k})_{m}]}{m!} + O \pfrac{\E[ (N_{\leq k})_M]}{M!} \\
 &= \sum_{m=0}^{M-1} (-1)^m \frac{ (\E N_{\leq k})^m}{m!}
 + O\pfrac{(\E N_{\leq k})^M}{M!} \\
 &\qquad + O_{k, M} \left( 
 \frac{B^{kM}}{\min\{c_1 + c^{1/2}, c'_1 + {c'}^{1/2}\}^{2/3}}\right) \\
 &= \expe^{-\E N_{\leq k}} 
 + O\pfrac{(\E N_{\leq k})^M}{M!} \\
 &\qquad + O_{k, M} \left( 
 \frac{B^{kM}}{\min\{c_1 + c^{1/2}, c'_1 + {c'}^{1/2}\}^{2/3}}\right).
\end{align*}
Since $\E N_{\leq k} = O_\eps(1)$ (by Lemma~\ref{lem:Nk-bound-small-k} again), we can choose $M = O_\eps(1)$ so that the first error term here is smaller than $\eps$. Note also $B = O(1)$ by hypothesis. Putting all this together, we have
\[
 \P(N = 0) = \expe^{-\E N} + O(\eps) + O_{\eps}\left(\min\{c_1 + c^{1/2}, c'_1 + {c'}^{1/2}\}^{-2/3} \right).
\]

Finally, note that the hypothesis $c_1 + c'_1 = o(n^{2/3})$ ensures that either
\[
 (c_1 + c^{1/2}) (c'_1 + {c'}^{1/2}) = o(n),
\]
in which case we are already done by Theorem~\ref{thm:main-ENk-bound}, or
\[
 \min\{c_1 + c^{1/2}, c'_1 + {c'}^{1/2}\} \gg n^{1/3}.
\]
Thus
\[
 \P(N = 0) = \expe^{-\E N} + o(1),
\]
as claimed.
\end{proof}

\end{document}
